\documentclass[11pt, a4paper]{article}

\usepackage[T1]{fontenc}
\usepackage[utf8]{inputenc}
\usepackage{tgpagella}
\usepackage{tgcursor}
\usepackage{microtype}
\usepackage[spanish, es-noshorthands]{babel}
\addto\captionsspanish{\renewcommand{\tablename}{Tabla}}

\usepackage{amsmath, amssymb}
\usepackage[table, xcdraw]{xcolor}
\usepackage{graphicx}
\usepackage{booktabs}
\usepackage{tabularx}
\usepackage[most]{tcolorbox}
\usepackage[margin=2.5cm]{geometry}
\usepackage{fancyhdr}

\usepackage{listings}
\lstdefinestyle{xmlstyle}{
    language=XML,
    basicstyle=\ttfamily\small,
    backgroundcolor=\color{gray!5},
    frame=single,
    rulecolor=\color{gray!40},
    keywordstyle=\color{blue}\bfseries,
    stringstyle=\color{red!70!black},
    commentstyle=\color{green!60!black},
    showstringspaces=false,
    breaklines=true, % Soluciona desbordamientos en bloques de código
    breakatwhitespace=true,
    morekeywords={robot, link, joint, origin, axis, parent, child, limit}
}
\lstset{style=xmlstyle}

\pagestyle{fancy}
\fancyhf{}
\setlength{\headheight}{16pt}
\lhead{\textbf{UCB Sede Tarija} | Ing. Mecatrónica}
\chead{Robótica (IMT-342) - Guía Técnica}
\rhead{Soporte de Implementación}
\lfoot{Prof: M.Sc. Ing. Bernardo Quiroga Turdera}
\rfoot{Página \thepage}
\renewcommand{\headrulewidth}{0.4pt}
\renewcommand{\footrulewidth}{0.4pt}

\definecolor{ucbblue}{RGB}{0, 51, 102}

\begin{document}

\begin{center}
    \Large\textbf{\textcolor{ucbblue}{Guía de Construcción XML/URDF para el ABB IRB 120}}\\[0.2cm]
    \normalsize \textbf{Objetivo:} Construir progresivamente un modelo URDF del manipulador y demostrar que su árbol cinemático es coherente con el modelo matemático D-H[cite: 5].
\end{center}

\vspace{0.4cm}
\begin{tcolorbox}[colback=red!5, colframe=red!75!black, title=\textbf{Concepto Crítico D-H vs. URDF {[cite: 5]}}]
    \textbf{Una fila de parámetros D-H no puede copiarse directamente en un elemento \texttt{<origin>} de URDF.} D-H restringe la construcción del marco a 4 parámetros; URDF utiliza una transformación posicional explícita (\texttt{xyz/rpy}) y define mecánicamente el eje de rotación local.
\end{tcolorbox}

\section{Estructura Mínima URDF (Árbol Cinemático)[cite: 5]}
La estructura de URDF usa la jerarquía \texttt{parent} (eslabón superior) y \texttt{child} (eslabón inferior).
\begin{lstlisting}
<robot name="irb120_student_model">
  <link name="base_link"/>
  <link name="link_1"/>

  <joint name="joint_1" type="revolute">
    <parent link="base_link"/>
    <child link="link_1"/>
    
    <!-- Pose relativa del origen de J1 desde la base -->
    <origin xyz="0 0 0" rpy="0 0 0"/>
    
    <!-- Eje de rotacion de J1 en coordenadas locales -->
    <axis xyz="0 0 1"/>
    
    <!-- Limites fisicos (Requerira datos del datasheet ABB) -->
    <limit lower="[VERIFY]" upper="[VERIFY]" effort="[VERIFY]" velocity="[VERIFY]"/>
  </joint>
</robot>
\end{lstlisting}

\section{Distinción de Componentes[cite: 5]}
\begin{itemize}
    \item \textbf{\texttt{type="revolute"}}: Junta giratoria con límites angulares finitos (usado para los 6 motores del IRB 120).
    \item \textbf{Atributo \texttt{origin}}: Define el desplazamiento fijo en el espacio mediante \texttt{xyz} (metros) y \texttt{rpy} (radianes).
    \item \textbf{Atributo \texttt{axis}}: Define el vector local de giro. Ej: \texttt{<axis xyz="0 1 0"/>} rota sobre el eje Y.
\end{itemize}

\section{Flujo de Trabajo Incremental[cite: 5]}
No introduzca las 6 juntas a la vez. Siga este flujo:
\begin{enumerate}
    \item Cree el esqueleto y \texttt{base\_link}. Use \texttt{check\_urdf} para validar el parseo XML.
    \item Añada \texttt{link\_1} y \texttt{joint\_1}. Evalúe en RViz. Verifique la dirección del giro.
    \item Repita agregando el siguiente eslabón cuidando la cadena padre-hijo.
\end{enumerate}

\section{Hoja de Definición de Juntas (Evidencia Entregable)[cite: 5]}
Complete esta tabla razonando la relación física entre los eslabones:

\begin{center}
    \small % Reduccion de fuente para evitar desbordamientos
    \renewcommand{\arraystretch}{1.5}
    \begin{tabularx}{\textwidth}{|c|c|c|X|X|c|c|}
    \hline
    \textbf{Joint} & \textbf{Parent} & \textbf{Child} & \textbf{\texttt{xyz}} & \textbf{\texttt{rpy}} & \textbf{\texttt{axis}} & \textbf{Límites} \\ \hline
    J1 & base\_link & link\_1 & & & & [VERIFY] \\ \hline
    J2 & link\_1 & link\_2 & & & & [VERIFY] \\ \hline
    J3 & link\_2 & link\_3 & & & & [VERIFY] \\ \hline
    J4 & link\_3 & link\_4 & & & & [VERIFY] \\ \hline
    J5 & link\_4 & link\_5 & & & & [VERIFY] \\ \hline
    J6 & link\_5 & link\_6 & & & & [VERIFY] \\ \hline
    \end{tabularx}
\end{center}

\section{Validación Cruzada FK $\leftrightarrow$ URDF[cite: 5]}
Elija dos configuraciones articulares $\mathbf{q}^{(A)}$ y $\mathbf{q}^{(B)}$. Para cada una:
\begin{enumerate}
    \item Calcule la FK matemáticamente (${}^0T_6$).
    \item Posicione el robot URDF en los mismos valores articulares usando el \textit{Joint State Publisher}.
    \item Compare la posición/orientación resultante de la brida.
\end{enumerate}

\begin{center}
    \begin{tabularx}{\textwidth}{|c|X|X|c|}
    \hline
    \textbf{Conf.} & \textbf{Posición FK calculada} & \textbf{Observación en URDF/RViz} & \textbf{¿Coinciden?} \\ \hline
    A & $p = [ \quad, \quad, \quad ]$ & $p = [ \quad, \quad, \quad ]$ & \\ \hline
    B & $p = [ \quad, \quad, \quad ]$ & $p = [ \quad, \quad, \quad ]$ & \\ \hline
    \end{tabularx}
\end{center}

\end{document}
